\documentclass[11pt]{article}
\usepackage[T1]{fontenc}
\usepackage{lmodern}
\usepackage[margin=1in]{geometry}
\usepackage{amsmath,amssymb,amsthm,booktabs,array}
\usepackage[colorlinks=true,urlcolor=blue,citecolor=blue,linkcolor=blue]{hyperref}
\newtheorem{theorem}{Theorem}
\newtheorem{lemma}{Lemma}
\newtheorem{corollary}{Corollary}
\newcommand{\val}{\nu_2}
\title{A restoring and divisor-sum proof of Bala's\newline
  machine-3626 conjecture of 11 September 2026,\newline
  OEIS A393210, Comments, revision 18:\newline
  extending Sequence Machine (2026) and Wolfram (2026)}
\author{Haobo Ma\\\small ChronoAI Pte Ltd, Singapore
  \and Wenlin Zhang\\\small National University of Singapore\\
  \small ChronoAI Pte Ltd, Singapore}
\date{4 November 2026}
\begin{document}
\maketitle
\begin{abstract}
Extending the pointwise formula recorded by Sequence Machine (18 April 2026)
and Wolfram's (2026) worst-case analysis, we prove Peter Bala's conjecture
dated 11 September 2026 in the Comments
of OEIS A393210, revision 18, also stated in A069517, Formula, revision 31.
For every positive integer $n$, the one-sided Turing machine 3626 halts after
$4\cdot2^{\val(n)}-3$ transitions, including its final boundary crossing.
A restoring block computation gives the recurrence $T_{k+1}=2T_k+3$;
the signed M\"obius--totient divisor sum in the conjecture has the same value.
The proof also gives exact timing statistics and the formulas for the two
comparison machines in the same comment. A general weighted divisor identity
gives the companion formulas in A069517 and A296306, including a correction
to the latter entry's stated value on powers of two.
The companion ordinary generating series recover Stephan's (2003) family.
\end{abstract}

\section{The machine and the conjecture}
The halting-time sequence A393210 was entered by Sean A. Irvine in February
2026. Bala's September comment proposes the divisor formula below
\cite{halting,mobius}. Wolfram's earlier analysis of identity-computing
machines displays $2^{\ell+1}-3$ in its closed-form worst-case table,
given for even lengths with analogous odd-length forms. Its named growth
class and identity-runtime profiles identify machine 3626 with this case
\cite{wolfram}.
The restoring argument below gives the time for each individual input and
identifies it with Bala's signed divisor sum.
The arithmetic ingredient is already explicit in A069517: Antti Karttunen's
corrected support formula of 19 November 2017 and Amiram Eldar's prime-power formula of
29 August 2023 describe its values.
We give a direct proof of that description and connect it to the machine's
complete computation.

Sequence Machine's archived A393210 record of 18 April 2026 already gives
the zero-based formula $a(k)=2\operatorname{xor}(k,k+1)-1$, with comparisons
against 10,000 terms \cite{sequence-machine}. For $n=k+1$, subtraction
flips the trailing $\val(n)$ zeros and the first one, so
$\operatorname{xor}(n-1,n)=2^{\val(n)+1}-1$. This recorded expression is
therefore the same pointwise closed form. It is a prior answer supported
by finite comparisons; the restoring trace and the divisor-sum equality
are proved below.

Number the tape cells $0,1,2,\ldots$, increasing to the left. Initially the
head is at cell $0$ in state $A$; cell $j$ holds the coefficient of $2^j$
in $n$. All remaining cells are zero. A move right from cell $0$ terminates
the computation and counts as one transition. There is no separate halt
state. The rules are
\[
\begin{array}{c|cc}
 &0&1\\\hline
A&(0,L,A)&(1,R,B)\\
B&(1,L,A)&(0,R,B)
\end{array}
\]
where each triple specifies the written symbol, direction and new state.
The tape is infinite to the left; there is no tape-length restriction.
Write $H(n)$ for the first boundary-crossing time, and put
\[
 m(n)=\sum_{e\mid n}(-1)^{e+1}\mu(e),
 \qquad C(n)=\sum_{d\mid n}m(d)^2\varphi(d).
\]
Here $\mu$ is the M\"obius function and $\varphi$ is Euler's totient.
Bala's conjecture is $H(n)=C(n)$ for every $n\geq1$.

\begin{theorem}\label{thm:main}
For every positive integer $n$, machine 3626 terminates, leaves every tape
symbol unchanged, and satisfies
\[
 H(n)=C(n)=4\cdot2^{\val(n)}-3.
\]
\end{theorem}

\section{A restoring computation}
Let $u_k=0^k1w$, where $w$ is an arbitrary finite string of higher bits,
with implicit zeros beyond it. Strings are read from the head toward the
left. Write $(S,u\mid r)$ for a configuration in which the scanned symbol
is the first symbol of $u$, the state is $S$, and $r$ lists the cells to
the right, nearest first. For a string $u$ define
\[
 E(u,r)=
 \begin{cases}
 \text{termination},&r=\varnothing,\\
 (B,bu\mid r'),&r=br'.
 \end{cases}
\]
This is the endpoint of moving right past the starting cell while restoring
that cell and the cells to its left.

\begin{lemma}\label{lem:block}
Starting in $(A,u_k\mid r)$, the machine first moves right past the starting
cell after exactly $T_k$ transitions, reaches $E(u_k,r)$, and restores every
tape symbol. The time is independent of $w$ and $r$, with
\[
 T_0=1,\qquad T_{k+1}=2T_k+3,\qquad T_k=4\cdot2^k-3.
\]
\end{lemma}
\begin{proof}
For $k=0$, the scanned symbol is $1$. Rule $A1$ moves right without changing
it, giving the assertion in one transition.

Assume the assertion for $k$, and write $u=u_k$. For the block $0u$, the
entire computation is
\[
\begin{aligned}
(A,0u\mid r)&\xrightarrow{1}(A,u\mid0r)
\xrightarrow{T_k}(B,0u\mid r)\\
&\xrightarrow{1}(A,u\mid1r)
\xrightarrow{T_k}(B,1u\mid r)
\xrightarrow{1}E(0u,r).
\end{aligned}
\]
The one-step arrows use $A0$, $B0$, and $B1$, respectively. Each long arrow
is the induction hypothesis with the indicated additional right-hand cell.
That cell prevents termination during either subcomputation. Each
subcomputation restores $u$ and first returns right only to the extra cell.
The extra cell changes from $0$ to $1$ and back to $0$. Thus all symbols are
restored, and only the last arrow moves right past the original starting
cell. This proves both the first-crossing property and the recurrence.
Finally $T_{k+1}+3=2(T_k+3)$ and $T_0+3=4$, giving the closed form.
\end{proof}

For $n>0$, its least-significant-bit-first expansion is uniquely
$0^{\val(n)}1w$. Apply Lemma~\ref{lem:block} with $r=\varnothing$.
This proves termination, restoration, and the asserted value of $H(n)$.
The excluded input $n=0$ consists entirely of zeros: rule $A0$ moves left
forever, so this input does not terminate.

\begin{corollary}[The comparison machines]\label{cor:machines}
Under the same initialization and boundary convention, the two other
identity-computing machines in Bala's A393210 comment have exact times
\[
 H_{3589}(n)=2\val(n)+1,\qquad
 H_{3197}(n)=2-(-1)^n\quad(n\geq1).
\]
Both restore the tape. These are A037227 and A010684, respectively
\cite{comparisons}.
\end{corollary}
\begin{proof}
The rule numbers use four base-eight digits, in the order $A1,A0,B1,B0$.
A digit $b$ specifies state $1+\lfloor b/4\rfloor$, symbol
$\lfloor b/2\rfloor\bmod2$, and direction right exactly when $b$ is odd
\cite{decoder}. For 3589 the rules are
\[
 A0\mapsto(0,L,A),\quad A1\mapsto(1,R,B),\quad
 B0\mapsto(0,R,B),\quad B1\mapsto(0,L,A).
\]
On $0^v1w$, there are $v$ left moves across zeros, one right move at the
first one, then $v$ right moves across the unchanged zeros. Rule $B1$
is never used. For 3197, $A0\mapsto(0,R,A)$ and
$A1\mapsto(1,L,B)$; either $B$ rule preserves its symbol and moves right.
An even input therefore exits in one move; an odd input makes one left
and two right moves. No tape symbol changes in either computation.
\end{proof}

\begin{corollary}[Timing distribution]\label{cor:timing}
For $N\geq1$, the number of inputs $n\leq N$ with
$H(n)=4\cdot2^k-3$ is
\[
 \left\lfloor\frac{N}{2^k}\right\rfloor-
 \left\lfloor\frac{N}{2^{k+1}}\right\rfloor\quad(k\geq0).
\]
Furthermore,
\begin{equation}\label{eq:time-sum}
 \sum_{n=1}^N H(n)=N+4\sum_{j\geq1}2^{j-1}
                    \left\lfloor\frac{N}{2^j}\right\rfloor .
\end{equation}
Among the $2^{\ell-1}$ inputs of binary length $\ell\geq1$, the mean
time is $2\ell-1$ and the maximum is $2^{\ell+1}-3$.
The maximum reproduces Wolfram's earlier closed form \cite{wolfram}; the
distribution and mean follow from the individual-input formula above.
\end{corollary}
\begin{proof}
Count multiples of $2^k$ that are not multiples of $2^{k+1}$.
The identity $2^v=1+\sum_{j=1}^v2^{j-1}$ gives
\eqref{eq:time-sum}. For $N=2^h-1$ this is
$(2h-3)2^h+3$, also valid at $h=0$. Subtract the expressions for
$h=\ell$ and $h=\ell-1$ to obtain the mean. The maximum occurs at the
unique length-$\ell$ input $n=2^{\ell-1}$.
\end{proof}

\section{The divisor sum and its companions}
Write $n=2^vq$ with $q$ odd. In the sum defining $m(n)$, only divisors
$e$ with at most one factor of $2$ contribute. If $v=0$, their contribution
is $\sum_{e\mid q}\mu(e)$. If $v\geq1$, pair each odd divisor $e$ with
$2e$: since $\mu(2e)=-\mu(e)$, each pair contributes $2\mu(e)$.
For $q>1$, the sum $\sum_{e\mid q}\mu(e)$ is zero by expansion of
$\prod_{p\mid q}(1-1)$; for $q=1$ it is one. Consequently
\begin{equation}\label{eq:support}
m(1)=1,\qquad m(2^j)=2\ (j\geq1),\qquad m(n)=0
\ \text{otherwise}.
\end{equation}
This also proves multiplicativity: two coprime arguments cannot both be
even; an odd argument greater than one forces both sides of the
multiplicativity identity to be zero. Thus
\[
C(n)=1+4\sum_{j=1}^{v}\varphi(2^j)
=1+4\sum_{j=1}^{v}2^{j-1}
=4\cdot2^v-3,
\]
including $v=0$ through the empty-sum convention. This completes the proof
of Theorem~\ref{thm:main}.

The other formulas in Bala's A069517 block follow from the same support
calculation, without further machine analysis.
\begin{corollary}[Arbitrary divisor weights]\label{cor:weights}
For $n\geq1$, $v=\val(n)$, every integer $r\geq1$, and every function
$w$ on the positive integers with complex values,
\begin{equation}\label{eq:weights}
 \sum_{d\mid n}m(d)^r w(d)=w(1)+2^r\sum_{j=1}^v w(2^j).
\end{equation}
\end{corollary}
\begin{proof}
Only $d=1$ and the positive powers of two contribute, by
\eqref{eq:support}. Their values are $1$ and $2$, respectively.
\end{proof}

\begin{corollary}\label{cor:mixed}
For integers $r\geq1$, $s\geq0$ and $n\geq1$, with $v=\val(n)$,
\[
\sum_{d\mid n}m(d)^r\varphi(d)^s=
\begin{cases}
1+2^rv,&s=0,\\[2pt]
1+2^r\dfrac{2^{sv}-1}{2^s-1},&s>0.
\end{cases}
\]
For every integer $k\geq1$,
$\sum_{d\mid n}(2m(d)^k-m(d)^{k+1})=1$.
\end{corollary}
\begin{proof}
Equation~\eqref{eq:support} reduces the first sum to
$1+2^r\sum_{j=1}^v2^{s(j-1)}$. Sum this geometric progression, or count its
$v$ terms when $s=0$. For the second sum, a divisor with $m(d)=0$ or $2$
contributes zero, and $d=1$ contributes one.
\end{proof}
In particular, the pairs $(r,s)=(1,1),(1,2),(2,2)$ give
\[
2^{v+1}-1,\qquad \frac{2\cdot4^v+1}{3},\qquad
\frac{4^{v+1}-1}{3},
\]
the entries A038712, A276391 and A296306, respectively
\cite{companions,ratio}. These are arithmetic identities attached to
Bala's conjecture, rather than further halting-time conjectures.

There is a useful distinction between the two roles of the power of two
in these formulas. The value is $m(2^j)=2$ for every $j\geq1$, not
$2^j$. The latter expression appears in Bala's 11 September 2026 comment
in A296306 and fails already at $j=2$, since $m(4)=2\ne4$.
The entry's displayed divisor identities nevertheless hold, with the
constant value $2$ used above.

\begin{corollary}[Signed transforms and indexing]\label{cor:signed}
For every integer $r\geq1$,
\begin{equation}\label{eq:signed}
 \sum_{d\mid n}m(d)^r\mu(d)=1-2^r\mathbf{1}_{2\mid n}.
\end{equation}
Using $w(d)=(-1)^{d+1}\mu(d)$ instead gives
\[
 \sum_{d\mid n}(-1)^{d+1}m(d)^r\mu(d)
 =1+2^r\mathbf{1}_{2\mid n}.
\]
This is the general formula in Bala's A010684 block of 22 September 2026.
Its case $r=1$, and also the expression
\[
 (-1)^{n+1}\sum_{d\mid n}m(d)^2\mu(d),
\]
equal $1$ for odd $n$ and $3$ for even $n$.
That block explicitly uses offset 1: in the entry's usual offset-0
notation these expressions are A010684$(n-1)$, whereas
$H_{3197}(n)$ is A010684$(n)$ \cite{comparisons}.
\end{corollary}
\begin{proof}
Substitute $w(d)=\mu(d)$ and $w(d)=(-1)^{d+1}\mu(d)$ into
\eqref{eq:weights}. Since $\mu(2)=-1$ and $\mu(2^j)=0$ for $j>1$,
only $j=1$ can contribute. The stated cases follow by parity.
\end{proof}

\begin{corollary}[A038712 and A276391]\label{cor:companions}
Write $B(n)=(2\cdot4^v+1)/3$ and $D(n)=2^{v+1}-1$.
Then $B(2n)=4B(n)-1$,
$B(2n-1)=1$, $D(2n)=2D(n)+1$, and $D(2n-1)=1$.
Their generating functions, binary digit sums and summatory identities
are given by the formulas below.
\end{corollary}
\begin{proof}
The recurrences follow from $\val(2n)=\val(n)+1$ and
$\val(2n-1)=0$. Also $D(n)=2(n\mathbin{\&}(-n))-1$, where the bitwise
expression denotes the largest power of two dividing $n$.
These are recorded formulas for A276391 and A038712
\cite{companions}. If $G(x)=\sum_{n\geq1}B(n)x^n$, separating odd
and even indices gives, for $|x|<1$,
\[
 G(x)=4G(x^2)+\frac{x}{1+x}
 =\sum_{j\geq0}\frac{4^j x^{2^j}}{1+x^{2^j}}.
\]
The series converges absolutely; iteration is justified since
$4^hG(x^{2^h})\to0$. This gives the generating-function identity
recorded by Ilya Gutkovskiy on 14 December 2020 in A276391. For the same $B$, let
$U(n)=\sum_{j\geq0}\epsilon_j4^j$ when $n=\sum_j\epsilon_j2^j$.
Subtracting the base-four digit strings of $n$ and $n-1$ yields
\[
 U(n)-U(n-1)=4^v-\sum_{j=0}^{v-1}4^j=B(n).
\]
Thus $\sum_{j=1}^n B(j)=U(n)$, the A000695 partial-sum relation
recorded in A276391 \cite{companions,moser}.
Also $H(n)=2D(n)-1$, so \eqref{eq:time-sum} gives the exact
summatory function of $D$ (A080277):
\[
 \sum_{n=1}^N D(n)=N+2\sum_{j\geq1}2^{j-1}\lfloor N/2^j\rfloor.
\]
Its value at $N=2^h-1$ is $(h-1)2^h+1$, and in general it is
$N\log_2N+O(N)$. Indeed, replace each floor by its argument; the
total error is bounded by $2\sum_{j\leq\log_2N}2^{j-1}=O(N)$.
This asymptotic appears as a conjecture in Brockhaus's 2003 comment in
A080277, whose 2013 update points to Bannister et al.\ for its proof
\cite{partial-sums}. The same exact sum gives
$S(N)=N+2S(\lfloor N/2\rfloor)$ and, for the binary digits of $N$,
\[
 S(N):=\sum_{n=1}^N D(n)=\sum_{i\geq0}\epsilon_i(i+1)2^i.
\]
This binary-digit formula also appears in the proof of Lemma~9 of
Bannister et al.\ \cite{partial-sums}.
Indeed, each digit's contribution to the floor sum occurs $i+1$ times.
The weights $(i+1)2^i$ are superincreasing, since the sum of the earlier
weights is $(i-1)2^i+1$ for $i\geq1$.
Their subset sums are therefore listed increasingly by $S(0),S(1),\ldots$;
this proves Bala's product identity of 2 February 2013 in A080277:
\[
 \prod_{i\geq0}(1+x^{(i+1)2^i})=\sum_{N\geq0}x^{S(N)}.
\]
\end{proof}

\begin{corollary}[Power sums and their series]\label{cor:series}
For real $a\geq0$, set $F_a(n)=\sum_{j=0}^{\val(n)}2^{aj}$.
Then
\begin{align*}
 \sum_{n\geq1}F_a(n)x^n
 &=\sum_{j\geq0}\frac{2^{aj}x^{2^j}}{1-x^{2^j}}
 &&(|x|<1),\\
 \sum_{n\geq1}\frac{F_a(n)}{n^s}
 &=\frac{\zeta(s)}{1-2^{a-s}}
 &&(\Re s>\max\{1,a\}).
\end{align*}
The cases $a=1,2,3$ are A038712, A296306, A318935, respectively.
Also $F_0(n)=\val(n)+1$
and $\sum_{d\mid n}m(d)=2F_0(n)-1=H_{3589}(n)$.
\end{corollary}
\begin{proof}
Write $F_a(n)=\sum_{2^j\mid n}2^{aj}$, interchange absolutely convergent
sums, and sum over positive multiples of $2^j$. For the last equality
use \eqref{eq:weights} with $r=1,w=1$.
\end{proof}
For positive integer $c$, the ordinary generating function with $c=2^a$
is Stephan's theorem of 2003, section~2, equation~(2.1), proved there by
the same coefficient-divisibility argument \cite{stephan}.
The formulation above also permits every real $a\geq0$.

\begin{corollary}[The four statements in A296306]\label{cor:ratio}
The divisor-power ratio in Conjecture~1 is an integer for every
nonnegative integer exponent. The value set, palindromic prefixes and
prefix sums in Conjectures~2--4 have the stated forms below.
The divisor factorization also holds at every prime and real exponent
$s\geq0$.
\end{corollary}
\begin{proof}
More generally, let $p$ be any prime, write $n=p^vq$ with $p\nmid q$,
and take any real $s\geq0$. Every divisor is uniquely $p^je$ with
$0\leq j\leq v$ and $e\mid q$. Hence
\begin{equation}\label{eq:ratio}
\frac{\sum_{d\mid n}d^s}{\sum_{e\mid n,\ p\nmid e}e^s}
=\sum_{j=0}^vp^{js}.
\end{equation}
The denominator is positive. Taking $p=2$ and integer $s\geq0$ makes
the right-hand side an integer, proving Conjecture~1 of A296306;
the factorization remains valid for every real $s\geq0$. For $s=2$, write
$R(n)=(4^{v+1}-1)/3$. As $v$ ranges over the nonnegative integers, this
gives exactly $(4^h-1)/3$, $h\geq1$, proving Conjecture~2.
For $1\leq j<2^h$, factoring out $2^{\val(j)}$ shows
$\val(2^h-j)=\val(j)$, so $R(2^h-j)=R(j)$; this proves Conjecture~3.
There are $2^{h-k-1}$ integers $1\leq j<2^h$ with $\val(j)=k$.
Therefore, for $h\geq1$,
\[
\sum_{j=1}^{2^h-1}R(j)
=\frac13\sum_{k=0}^{h-1}2^{h-k-1}(4^{k+1}-1)
=\frac{(2^h-1)(2^{h+1}-1)}3,
\]
which is A006095$(h+1)$ and proves Conjecture~4 \cite{gaussian}.
The entry was created by Ivan N. Ianakiev on 10 December 2017.
The existing record credits David Radcliffe's closed formula
(11 December 2017) and M. F. Hasler's explanations of Conjectures~1 and~3
(31 January 2020) \cite{ratio}; the deductions here use those same
elementary factorizations.
\end{proof}

\begin{corollary}[The cubic companion A318935]\label{cor:cubic}
For $p=2$, $s=3$, \eqref{eq:ratio} is
\[
 Q(n)=\frac{\sum_{d\mid n}d^3}{\sum_{e\mid n,\ e\text{ odd}}e^3}
     =\frac{8^{v+1}-1}{7}.
\]
This is A318935, introduced by N. J. A. Sloane on 14 September 2018
as the sum of cubes of the powers of two dividing $n$; Chai Wah Wu gave
its valuation formula that day. Amiram Eldar recorded the equivalent
cubic divisor ratio on 29 October 2026 \cite{cubic}.
The underlying cubic power sum already occurs as $T$ on page~102 of
Lamont's 1982 paper on Cayley integers \cite{lamont}.
It is also a direct weighted consequence:
\[
 Q(n)=1+4\sum_{\substack{d\mid n\\d>1}}m(d)\varphi(d)^3.
\]
\end{corollary}
\begin{proof}
The first equality is \eqref{eq:ratio}. In the weighted equality, the
sum is $2\sum_{j=0}^{v-1}8^j$. The same divisor
factorization gives the formula at every prime in \eqref{eq:ratio}.
\end{proof}

\section{Coverage and mathematical verification}
\begin{center}
\small
\begin{tabular}{@{}>{\raggedright\arraybackslash}p{0.39\linewidth}>{\raggedright\arraybackslash}p{0.55\linewidth}@{}}
\toprule
Posing item & Coverage\\\midrule
A393210, Comments rev.~18; A069517, Formula rev.~31 & Proved: one conjecture with two locators; Theorem~\ref{thm:main}, all $n>0$.\\
Sequence Machine, A393210 (18 Apr 2026) & Prior pointwise answer cited; all-input restoring and divisor-sum proof in Theorem~\ref{thm:main}.\\
A393210, identity-function and comparison comments & Reproved; prior results Wolfram (2026): restoration, Lemma~\ref{lem:block}; both other machines, Corollary~\ref{cor:machines}.\\
Wolfram (2026), worst-case timing of machine 3626 & Reproved; prior closed form Wolfram (2026): $\max H=2^{\ell+1}-3$, Corollary~\ref{cor:timing}.\\
A069517, four companion formulas (11 Sep 2026) & Proved: Corollaries~\ref{cor:weights}--\ref{cor:mixed}; support formulas of 2017 and 2023 cited.\\
A296306, comment $m(2^j)=2^j$ & Refuted-and-corrected: $m(4)=2$; the correct value is $2$ for every $j>0$.\\
A296306, Conjectures 1 and 3 & Reproved; prior proof Hasler (2020): Corollary~\ref{cor:ratio}.\\
A296306, Conjectures 2 and 4 & Proved from Radcliffe's prior formula (2017): Corollary~\ref{cor:ratio}.\\
A010684, Bala block (22 Sep 2026) & Proved with its declared reindexing: Corollary~\ref{cor:signed}.\\
A038712/A276391, valuation, recurrence, series and A000695 links & Proved: Corollaries~\ref{cor:companions}--\ref{cor:series}.\\
Divisor-power ordinary generating series & Reproved for integer $c=2^a$; prior proof Stephan (2003), eq.~(2.1); real-$a$ formulation in Corollary~\ref{cor:series}.\\
A080277, Bala's two divisor formulas; floor sum, binary product and asymptotic & Proved; asymptotic and binary-digit formula reproved, prior proof Bannister et al.\ (2013): Corollary~\ref{cor:companions}.\\
A318935, cubic ratio (29 Oct 2026) & Proved; earlier cubic sum Lamont (1982), valuation Wu (2018), ratio Eldar (2026): Corollary~\ref{cor:cubic}.\\
A000695, Rascoe's tree conjecture (2022); other combinatorial models; A038712 machine 924 & Not covered; different objects or transition rules.\\
\bottomrule
\end{tabular}
\end{center}

The pinned \href{https://github.com/the-omega-institute/trureturing/blob/3d88df0f99ef3c8bfb034fa52be669b13746a844/D5/S0/Computability/BinaryMachine/Machine3626.lean}{\texttt{Machine3626.lean}}
contains these declarations in the namespace
\begin{center}
\small\texttt{D5.S0.Computability.BinaryMachine.Machine3626}.
\end{center}
\begin{itemize}
\item \texttt{block\_first\_crossing}: for every $k$ and arbitrary finite higher
bits, a trace of length $4\cdot2^k-3$ from $0^k1w$ at the actual boundary
to its first crossing; the right-context list is empty.
\item \texttt{positive\_input\_first\_crossing}: for every $n>0$, existence
of $k,w$ with $n$'s binary list equal to $0^k1w$, and the corresponding
trace. The zero-block length is the $k=\val(n)$ identified above.
\item \texttt{result}: for every $n>0$, existence of a terminating trace
whose natural length, cast to the integers, equals the double signed
M\"obius--totient divisor sum in Theorem~\ref{thm:main}.
\end{itemize}
The deterministic step rules and the trace constructors allow no continuation
after crossing. These three declarations state trace and timing properties;
the endpoint \texttt{none} does not itself contain a final tape.
Lemma~\ref{lem:block} supplies the written restoration argument for all
contexts. Corollaries~\ref{cor:machines}--\ref{cor:cubic}, including the
comparison machines and arithmetic generalizations, have written proofs here.
The finite checks in \texttt{verify.py} compare all three literal transition
tables on $1\leq n\leq65536$, evaluate divisor identities on
$1\leq n\leq4096$, test restoring blocks with short contexts, and check the
timing statistics, arithmetic recurrences and indexing identities.
These finite comparisons supplement the all-input proofs.

\begin{thebibliography}{14}
\bibitem{halting} S. A. Irvine, OEIS A393210, entered 5 February 2026;
P. Bala, conjecture comment, 11 September 2026, Comments revision 18.
\url{https://oeis.org/A393210}.
\bibitem{mobius} B. Cloitre, OEIS A069517, entered 16 April 2002;
A. Karttunen, corrected support formula, 19 November 2017;
A. Eldar, prime-power formula, 29 August 2023; P. Bala, formulas,
11 September 2026, Formula revision 31. \url{https://oeis.org/A069517}.
\bibitem{wolfram} S. Wolfram, \emph{P vs. NP and the Difficulty of
Computation: A Ruliological Approach}, 30 January 2026, sections
``The Basic Setup,'' ``Runtimes in $s=2$, $k=2$ Machines''
(closed-form runtime table), and ``Computing the Same Functions at
Different Speeds'' (machines 3197, 3589, 3626).
\href{https://writings.stephenwolfram.com/2026/01/p-vs-np-and-the-difficulty-of-computation-a-ruliological-approach/}{Stephen Wolfram Writings}.
\bibitem{companions} The OEIS Foundation Inc., entries
\href{https://oeis.org/A038712}{A038712} and
\href{https://oeis.org/A276391}{A276391}, formulas in terms of the
exponent of $2$ dividing the argument.
\bibitem{ratio} I. N. Ianakiev, OEIS A296306, 10 December 2017;
D. Radcliffe, closed formula, 11 December 2017; M. F. Hasler,
comments, 31 January 2020. \url{https://oeis.org/A296306}.
\bibitem{gaussian} N. J. A. Sloane, OEIS A006095, formula
$a(h)=(2^h-1)(2^{h-1}-1)/3$. \url{https://oeis.org/A006095}.
\bibitem{comparisons} The OEIS Foundation Inc.,
\href{https://oeis.org/A037227}{A037227} (machine 3589) and
\href{https://oeis.org/A010684}{A010684} (machine 3197);
P. Bala, arithmetic formulas in A010684, 22 September 2026.
\bibitem{decoder} Wolfram Institute, \emph{TuringMachine},
\texttt{from\_number}; one-sided function
\texttt{OneSidedTuringMachine}\allowbreak\texttt{FunctionNative},
revision \texttt{d23b7bcb741c711ea7b87}\allowbreak\texttt{a76391da1ac150c3bba}.
\href{https://github.com/WolframInstitute/TuringMachine/blob/d23b7bcb741c711ea7b87a76391da1ac150c3bba/TuringMachine/Libs/ndtm_search/src/models.rs}{Rule-number decoder};
\href{https://github.com/WolframInstitute/TuringMachine/blob/d23b7bcb741c711ea7b87a76391da1ac150c3bba/TuringMachine/Kernel/Functions.wl}{one-sided initialization and counting}.
\bibitem{moser} The OEIS Foundation Inc., A000695, binary-to-base-four
digit interpretation. \url{https://oeis.org/A000695}.
\bibitem{cubic} N. J. A. Sloane, OEIS A318935, 14 September 2018;
C. W. Wu, valuation formula, same day;
A. Eldar, cubic divisor-ratio formula, 29 October 2026.
\url{https://oeis.org/A318935}.
\bibitem{partial-sums} The OEIS Foundation Inc., A080277,
formulas and the 2013 pointer to the proof of its asymptotic.
\url{https://oeis.org/A080277}. See also M. J. Bannister,
Z. Cheng, W. E. Devanny, and D. Eppstein,
\emph{Superpatterns and universal point sets},
\emph{J. Graph Algorithms Appl.} 18(2) (2014), 177--209,
section~2.3, Lemma~9, p.~182.
\href{https://doi.org/10.7155/jgaa.00318}{doi:10.7155/jgaa.00318};
arXiv:1308.0403, first posted 2 August 2013.
\url{https://arxiv.org/abs/1308.0403}.
\bibitem{stephan} R. Stephan, \emph{Divide-and-conquer generating functions.
Part I. Elementary sequences}, 2 July 2003, section~2, equation~(2.1),
pp.~3--4. \url{https://arxiv.org/abs/math/0307027}.
\bibitem{lamont} P. J. C. Lamont, \emph{The number of Cayley integers of
given norm}, \emph{Proc. Edinburgh Math. Soc.} 25(1) (1982), 101--103,
p.~102, the quantity $T$.
\href{https://doi.org/10.1017/S001309150000420X}{doi:10.1017/S001309150000420X}.
\bibitem{sequence-machine} Sequence Machine, A393210 formula record,
18 April 2026, revision \texttt{adf7a8346dc7}.
\href{https://github.com/jonmaiga/sequence-machine-data/blob/adf7a8346dc74ed9fec9462652611d740e33ec0b/oeis/A393/A393210.programs.json}{Archived formula record};
\href{https://sequencedb.net/s/A393210}{current sequence entry}.
\end{thebibliography}
\noindent\small\textcopyright\ 2026 Haobo Ma and Wenlin Zhang.
Licensed under Creative Commons Attribution 4.0 International (CC BY 4.0).
\end{document}
